\section{Method}
\label{headings}

\begin{figure*}[!t]
    \centering
    \includegraphics[width=1.0\linewidth]{images/structure.jpg}
    \vspace{-.3in}
    \caption{\textbf{The DreamNav architecture.} DreamNav is made up of two major components: the \emph{coarse-estimation stage} (Section~\ref{Coarse-estimation}), and the \emph{imagination-refined stage} (Section~\ref{sec:refine}). The first component uses \emph{rotation-aware matching} to map the source and target images into a calibrated two-view geometry, and then uses learned priors, metric scale cues and a heading network to create a metric pose hypothesis $\mathbf{p}_{\text{coarse}}$. The imagination-refined stage renders the view implied by this hypothesis with a ground-plane homography, re-matches it against the real observation tile by tile, then finds the refined pose $\mathbf{p}^\star$ using reliability-weighted fusion.}
    \label{fig:framework}
    \vspace{-.1in}
\end{figure*}

\textbf{Overview.} DreamNav uses a combination of learned priors and calibrated two-view geometry to map a source--target image pair to a metric pose change. This mapping comprises two stages (Figure~\ref{fig:framework}): (1) a coarse-estimation stage, which distils sparse correspondences and learned priors into an initial pose hypothesis, and (2) an imagination-refined stage that renders the view implied by this hypothesis and corrects the pose from its discrepancy with the real observation. Note that the imagined view is rendered by a closed-form geometric warp, which radically differs from generative imagination but keeps the geometry exact on the ground plane (see Section~\ref{sec:refine}).
% \textbf{概述。} DreamNav 将学习到的先验与已标定的两视图几何结合起来，把源图--目标图图像对映射为度量的位姿变化。该映射包含两个阶段（图~\ref{fig:framework}）：（1）粗估计阶段，把稀疏对应点和学习先验提炼为初始位姿假设；（2）想象精修阶段，渲染该假设所蕴含的视图，并根据它与真实观测的差异修正位姿。注意，想象视图由闭式的几何变换渲染得到，这与生成式想象有根本区别，但能保证地面上的几何是精确的。

\subsection{Problem Formulation}\label{PoseRep}
Our goal is to estimate the pose change $\mathbf{p}=(\Delta x,\Delta y,\Delta z,\Delta\phi,\Delta\theta)$ from a source image $S$ to a target image $G$, where $(\Delta x,\Delta y,\Delta z)$ is the metric displacement in the East--North--Up frame and $\Delta\phi$ and $\Delta\theta$ are the changes in camera tilt and heading. We consider image pairs with zero camera roll and a shared pinhole calibration $\mathbf{K}$, which is estimated once from training pairs. Considering that the source and target cameras are tilted by $\phi_S$ and $\phi_G$ from nadir, with $\Delta\phi=\phi_G-\phi_S$, we can write the relative rotation $\mathbf{R}$ of the rigid motion $\mathbf{X}_G=\mathbf{R}\mathbf{X}_S+\mathbf{t}$ between the two camera frames as:
\begin{equation}
\mathbf{R}=\mathbf{A}(\phi_G)^{\top}\,\mathbf{R}_z(\Delta\theta)\,\mathbf{A}(\phi_S),
\label{eq:rot}
\end{equation}
where $\mathbf{A}(\phi)$ is the rotation from a camera tilted by $\phi$ to its gravity-aligned frame and $\mathbf{R}_z(\cdot)$ is a rotation about the vertical axis. Under this factorization, we reduce the problem to estimating a relative pose $(\Delta\theta,\phi_S,\phi_G,\mathbf{u},d)$ consisting of three angles, a unit translation direction $\mathbf{u}\in\mathbb{S}^2$ expressed in the gravity-aligned frame of the source camera, and a metric distance $d>0$. Given the absolute heading $\psi_S$ of the source camera, the displacement in the world frame can be straightforwardly obtained from the levelled displacement $(\ell_x,\ell_y,\Delta z)^\top=d\,\mathbf{u}$ as
\begin{equation}
\begin{pmatrix}\Delta x\\ \Delta y\end{pmatrix}=
\begin{pmatrix}\cos\psi_S & \sin\psi_S\\ -\sin\psi_S & \cos\psi_S\end{pmatrix}
\begin{pmatrix}\ell_x\\ \ell_y\end{pmatrix}.
\label{eq:world}
\end{equation}
Given two calibrated images, it is possible to estimate the three angles and $\mathbf{u}$ exactly---in theory---and the displacement only up to its scale $d$~\citep{hartley2003multiple} and the absolute heading $\psi_S$. We should thus focus learning on these two quantities and on the reliability of each estimate, making the remaining pose parameters tightly constrained by geometry.
% \subsection{问题定义} 目标是估计从源图 $S$ 到目标图 $G$ 的位姿变化 $\mathbf{p}=(\Delta x,\Delta y,\Delta z,\Delta\phi,\Delta\theta)$：东-北-天坐标系下的度量位移、俯仰变化和航向变化。相机无横滚、共享同一针孔内参 $\mathbf{K}$（由训练图像对一次性标定）。相对旋转可分解为式~\eqref{eq:rot}；于是问题化为估计相对位姿 $(\Delta\theta,\phi_S,\phi_G,\mathbf{u},d)$。已知源相机绝对航向 $\psi_S$ 时，世界系位移由式~\eqref{eq:world} 得到。两张已标定图像理论上可以精确给出三个角度和方向 $\mathbf{u}$，但位移的尺度 $d$ 和绝对航向 $\psi_S$ 不可观；因此学习应集中在这两个量以及各估计的可靠性上，其余参数由几何严格约束。

\subsection{Coarse-Estimation Stage}\label{Coarse-estimation}
Given the image pair $S$ and $G$, the coarse stage uses sparse correspondences and learned priors to produce an initial estimate $\mathbf{p}_{\text{coarse}}$.
% \subsection{粗估计阶段} 给定图像对 $S$ 和 $G$，粗估计阶段利用稀疏对应点和学习先验得到初始估计 $\mathbf{p}_{\text{coarse}}$。

\noindent\textbf{Rotation-Aware Multi-Scale Matching.}
Directly applying a pretrained SuperPoint--SuperGlue matcher~\citep{detone2018superpoint,sarlin2020superglue} in the context of UAV image pairs is impractical, because neither network is invariant to in-plane rotation while the heading change spans the full circle. To remedy this problem, we propose to use a rotation-aware, coarse-to-fine inference scheme around the frozen matcher. To find the rotation, we apply the matcher to a low-resolution copy of the pair, rotating $G$ by $k\cdot90^\circ$ for $k\in\{0,1,2,3\}$, yielding four match sets $\mathcal{M}^{(k)}=\{(\mathbf{x}_i,\mathbf{x}'_i,w_i)\}$ with confidences $w_i$. We select a single rotation by maximum total confidence: $k^\star=\operatorname*{argmax}_{k}\sum_{i\in\mathcal{M}^{(k)}}w_i$. After selecting $k^\star$, the pair is matched again at a higher resolution with a larger keypoint budget, and all keypoints are mapped back to the unrotated image at the original resolution. Coarse matches are thus responsible for the robustness of the pose estimate while fine matches enhance its accuracy. To prevent imprecise coarse keypoints from degrading well-matched pairs, we add the coarse matches to the pose estimation only if the fine matches alone yield too few inliers.
% \noindent\textbf{旋转感知的多尺度匹配。} 直接把预训练的 SuperPoint--SuperGlue 用于无人机图像对并不可行：两个网络都不具备面内旋转不变性，而航向变化覆盖整个圆周。为此，我们在冻结的匹配器外面套一层旋转感知、由粗到精的推理方案：先在低分辨率下把 $G$ 旋转 $k\cdot90^\circ$ 各匹配一次，按置信度总和选出旋转 $k^\star$；再以更高分辨率、更多关键点重新匹配，并把关键点映射回未旋转的原始分辨率图像。粗匹配负责鲁棒性，精匹配负责精度；仅当精匹配内点过少时才把粗匹配加入位姿解算。

\noindent\textbf{Relative Pose from Two-View Geometry.}
Once the matches are established, the geometric estimation proceeds, using robust essential-matrix fitting with MAGSAC++~\citep{barath2020magsacpp} followed by a nonlinear refinement of the Sampson error~\citep{hartley2003multiple}. In resolving the four-fold decomposition of the essential matrix, some regularities of UAV imagery can be leveraged: triangulated points lie in front of both cameras, the camera roll is zero, and the tilt stays within a limited range. Since the roll is zero, each estimated $\mathbf{R}$ corresponds to a triplet $(\Delta\theta,\phi_S,\phi_G)$ obtained in closed form from the factorization in Eq.~\eqref{eq:rot}. This also lets us express the translation direction in the gravity-aligned frame of the source camera, $\mathbf{u}=-\mathbf{A}(\phi_S)\mathbf{R}^\top\mathbf{t}/\lVert\mathbf{t}\rVert_2$, whose vertical component gives the direction of $\Delta z$.
% \noindent\textbf{由两视图几何求相对位姿。} 匹配建立后，用 MAGSAC++ 鲁棒拟合本质矩阵，再对 Sampson 误差做非线性精修。在本质矩阵的四重分解中利用无人机图像的规律选解：三角化点位于两相机前方、横滚为零、俯仰角范围有限。由于横滚为零，每个 $\mathbf{R}$ 由式~\eqref{eq:rot} 以闭式对应一组 $(\Delta\theta,\phi_S,\phi_G)$；据此还可把平移方向表达在源相机的重力对齐坐标系中，其竖直分量给出 $\Delta z$ 的方向。

\noindent\textbf{MAP Pose with Learned Priors.}\label{sec:head}
Two-view pose estimation is ambiguous in pairs that exhibit a near-planar scene and a short baseline, as a pitch rotation then resembles a forward translation. Such ambiguity can be resolved by injecting learned priors on the heading change and on the tilt of each camera when the matches are uninformative. To this end, an MLP over pooled statistics of a frozen MASt3R model~\citep{leroy2024mast3r} predicts the priors $(\overline{\Delta\theta},\bar\phi_S,\bar\phi_G)$ along with a direct regression of every pose parameter. The total error over the $N$ inlier matches is
\begin{equation}
E=\sum_{i=1}^{N}\rho\!\left(\frac{e_i}{\sigma_e}\right)
+\frac{(\phi_S-\bar\phi_S)^2+(\phi_G-\bar\phi_G)^2}{\sigma_\phi^2}
+\frac{(\Delta\theta\ominus\overline{\Delta\theta})^2}{\sigma_\theta^2},
\label{eq:map}
\end{equation}
where $e_i$ is the Sampson error of match $i$ in pixels, $\ominus$ denotes the wrapped angular difference, and $\rho$ is the pseudo-Huber cost. This nonlinear least-squares cost is iteratively minimized from both the essential-matrix solution and the prior pose using a trust-region solver, and the better-supported solution is kept. The data term dominates if the matches tightly constrain the pose. The priors take over as the geometry degenerates, which pins the tilts and thereby stabilizes the translation direction.
% \noindent\textbf{带学习先验的 MAP 位姿。} 场景近似平面且基线较短时，两视图位姿估计存在歧义：俯仰旋转与前向平移相似。在匹配信息不足时，可以注入关于航向变化和各相机俯仰角的学习先验来消除歧义。为此，一个作用在冻结 MASt3R 输出统计量上的 MLP 预测先验 $(\overline{\Delta\theta},\bar\phi_S,\bar\phi_G)$，并同时直接回归每个位姿参数。式~\eqref{eq:map} 中 $e_i$ 是以像素计的 Sampson 误差，$\ominus$ 是取模后的角度差，$\rho$ 是 pseudo-Huber 代价。该非线性最小二乘问题分别从本质矩阵解和先验位姿出发，用信赖域方法迭代求解，保留支持度更高的解。匹配约束充分时数据项占主导；几何退化时先验接管，把俯仰角固定住，平移方向也随之稳定。

\noindent\textbf{Pose-Guided Zoom Re-Matching.}
Matching entire images is necessary to establish correspondences when nothing is known about the relative pose. However, refining the pose is an inherently local operation, which, we show, can substantially benefit from concentrating keypoints and resolution on the co-visible region. Given the rotation-only homography $\mathbf{K}\mathbf{R}^\top\mathbf{K}^{-1}$ of the MAP pose, we warp $G$ into the orientation of $S$ and predict the window in which the two views overlap. For every pair, we first locate this window in both images, and then crop and resample it from the original resolution to the input size of the matcher before matching again. The zoomed matches are not pooled with those of the first pass as we found the pose solved from the zoomed matches alone to be more accurate. We denote the two poses re-estimated from them as the zoom-pass solutions.
% \noindent\textbf{位姿引导的放大重匹配。} 在对相对位姿一无所知时，必须对整幅图像做匹配。但位姿精修本质上是局部操作，把关键点和分辨率集中到共视区域能带来明显收益。利用 MAP 位姿对应的纯旋转单应 $\mathbf{K}\mathbf{R}^\top\mathbf{K}^{-1}$，把 $G$ 变换到 $S$ 的朝向并预测两视图的重叠窗口；随后在两幅图中定位该窗口，从原始分辨率裁剪并重采样到匹配器的输入尺寸后再次匹配。放大后的匹配不与第一遍匹配合并（单独用放大匹配解出的位姿更准），由它们重新估计的两个位姿记为 zoom-pass 解。

\noindent\textbf{Metric Scale Recovery.}
The translation $\mathbf{t}$ is recovered up to an unknown scale factor. The metric distance $d$ is impossible to recover with a purely geometric approach. We thus use three complementary metric cues: 1)~The prior MLP regresses the distance from the MASt3R statistics and thus provides an estimate $d_{\text{reg}}$ without relying on any match. 2)~The baseline-to-depth ratio $s$ of the triangulated inliers is combined with the metric scene depth $D$ of MASt3R into an estimate $d_{\text{depth}}=c\,s\,D$ that makes no assumption on the flight altitude. 3)~We estimate $d_{\text{alt}}=H_0/\tilde h$ by comparing the height $\tilde h$ of the camera above the triangulated ground, in baseline units, against the typical flight altitude $H_0$. The third cue therefore encodes the flight altitude of the training data in a single constant, requiring recalibration at other altitudes.
% \noindent\textbf{度量尺度恢复。} 平移 $\mathbf{t}$ 只能恢复到相差一个未知尺度因子，纯几何方法无法得到度量距离 $d$。因此我们使用三条互补的度量线索：1）先验 MLP 由 MASt3R 统计量直接回归距离，得到不依赖任何匹配的估计 $d_{\text{reg}}$；2）三角化内点的“基线/深度”比 $s$ 与 MASt3R 的度量场景深度 $D$ 组合为 $d_{\text{depth}}=c\,s\,D$（$c$ 在训练集上标定），不对飞行高度做任何假设；3）把相机到三角化地面的高度 $\tilde h$（以基线为单位）与典型飞行高度 $H_0$ 比较，得到 $d_{\text{alt}}=H_0/\tilde h$。第三条线索把训练数据的飞行高度编码成了一个常数，换到其他高度需要重新标定。

\noindent\textbf{Reliability-Weighted Candidate Fusion.}
Committing to a single estimator for all image pairs would be brittle. We instead express each pose parameter as a reliability-weighted blend of candidate estimates followed by a learned residual:
\begin{equation}
\hat y=\bar y+r\bigl(\mathbf{f},\{c_k-\bar y\}_k\bigr),\qquad
\bar y=\sum_{k}\frac{\hat e_k^{-2}}{\sum_{j}\hat e_j^{-2}}\,c_k,
\label{eq:blend}
\end{equation}
where $k$ indexes over the candidates $c_k$ of a pose parameter $y\in\{\Delta\theta,\Delta\phi,\log d,\Delta z,\ell_x,\ell_y\}$, $\hat e_k$ is the absolute error of $c_k$ predicted from pair-level features $\mathbf{f}$ such as inlier counts, residuals and the agreement between candidates, and $\bar y$ and $r$ are the blended estimate and the residual regressor, respectively. The candidates consist of the essential-matrix and MAP solutions of both matching passes and the direct regressions from the MASt3R statistics, each with a different failure mode and whose displacement components are given as $\hat d\,\mathbf{u}$ with the fused distance $\hat d$. For each candidate, a GBDT~\citep{ke2017lightgbm} learns $\hat e_k$ by fitting the square root of its absolute error (a compromise between typical and tail errors). We use the $\ell_1$ loss for the residual regressor. For the heading change, the blend $\bar y$ and the residual are computed on the unit circle. This makes the fusion robust to failed estimators: candidates for which the model predicts a large error will have a smaller effect on the blend. During this fusion, only the gradient-boosted regressors are trained, whereas the weights of the matcher and of MASt3R are kept frozen. As each regressor never sees the scenes it predicts (five-fold cross-fitting grouped by scene), its outputs on training pairs follow the same error distribution as at test time.
% \noindent\textbf{按可靠性加权的候选融合。} 对所有图像对只依赖单一估计器是脆弱的。我们把每个位姿参数表示为候选估计按可靠性加权的混合，再加一个学习到的残差（式~\eqref{eq:blend}）：$c_k$ 是参数 $y$ 的第 $k$ 个候选，$\hat e_k$ 是由图像对级特征 $\mathbf{f}$（内点数、残差、候选之间的一致性等）预测的该候选的绝对误差，$\bar y$ 和 $r$ 分别是混合估计和残差回归器。候选包括两遍匹配的本质矩阵解与 MAP 解，以及由 MASt3R 统计量直接回归的结果；位移分量由 $\hat d\,\mathbf{u}$ 构成。每个候选的 $\hat e_k$ 由一个 GBDT 拟合其绝对误差的平方根得到；残差回归器使用 $\ell_1$ 损失；航向变化的混合与残差在单位圆上计算。这样，被预测为误差大的候选对混合结果的影响就小。融合阶段只训练梯度提升回归器，匹配器和 MASt3R 的权重保持冻结。每个回归器从未见过它所预测的场景（按场景分组的五折交叉拟合），因此它在训练对上的输出与测试时服从相同的误差分布。

\noindent\textbf{Absolute Heading.}
Converting the levelled displacement to the world frame requires extracting the absolute heading $\psi_S$ from single-image cues such as the direction of cast shadows. Our heading network, a ConvNeXt~\citep{liu2022convnet}, outputs a von Mises distribution over the heading of each image, given by a mean direction $\mu$ and a concentration $\kappa$. Given the heading labels of the training images, we minimize the negative log-likelihood of this distribution:
\begin{equation}
\mathcal{L}_{\psi}=-\kappa\cos(\psi-\mu)+\log I_0(\kappa),
\label{eq:vm}
\end{equation}
where $I_0$ is the modified Bessel function of the first kind of order zero. For the source camera, we express the joint estimate over the two single-image predictions as
\begin{equation}
\hat\psi_S=\arg\bigl(\kappa_S\,e^{\mathrm{i}\mu_S}+\kappa_G\,e^{\mathrm{i}(\mu_G-\Delta\hat\theta)}\bigr),
\label{eq:absheading}
\end{equation}
where $\Delta\hat\theta$ denotes the fused heading change. This is simply computed by rotating the prediction of the target image into the source camera and summing the natural parameters of the two distributions. After fusion, we rotate $(\hat\ell_x,\hat\ell_y)$ by $\hat\psi_S$ with Eq.~\eqref{eq:world} and recover $\mathbf{p}_{\text{coarse}}=(\Delta\hat x,\Delta\hat y,\Delta\hat z,\Delta\hat\phi,\Delta\hat\theta)$.
% \noindent\textbf{绝对航向。} 把水平位移转换到世界坐标系，需要从单张图像的线索（例如阴影方向）中提取绝对航向 $\psi_S$。航向网络（ConvNeXt）对每张图像输出航向的 von Mises 分布（均值方向 $\mu$、集中度 $\kappa$），用训练图像的航向标签最小化负对数似然（式~\eqref{eq:vm}）。对源相机，把两张图各自的预测联合起来（式~\eqref{eq:absheading}）：将目标图的预测用融合后的 $\Delta\hat\theta$ 旋转到源相机，再把两个分布的自然参数相加。融合之后，用 $\hat\psi_S$ 按式~\eqref{eq:world} 旋转 $(\hat\ell_x,\hat\ell_y)$，得到 $\mathbf{p}_{\text{coarse}}$。

\subsection{Imagination-Refined Stage}
\label{sec:refine}
Intuitively, if a pose hypothesis is correct, we can use it to imagine one view from the other, and the imagined view will coincide with the real observation. The key refinement signal of this stage comes from the task of \emph{novel view synthesis}: given one view of a scene and a pose hypothesis, synthesize the image of the scene seen from the other camera. This is a particularly demanding task in our setting because a displacement of one metre shifts the ground by only about 16 pixels at the native resolution---roughly two pixels at the 512-pixel resolution of latent diffusion models---which leaves little margin for the spatial inaccuracy of generated images. We propose to tackle the limited fidelity of generative imagination with a generator-free design.
% \subsection{想象精修阶段} 直观上，如果位姿假设是正确的，就可以用它从一个视图想象出另一个视图，且想象视图会与真实观测重合。本阶段的精修信号来自新视角合成任务。这在我们的设定下要求很高：一米的位移在原始分辨率下只让地面移动约 16 个像素，在潜扩散模型的 512 像素分辨率下只有约 2 个像素，几乎没有给生成图像的空间误差留下余量。因此我们用“无生成器”的设计来解决生成式想象保真度不足的问题。

\noindent\textbf{Pose Hypotheses.}\label{heuristic}
We are given the pose estimates of the coarse stage, namely the fused estimate $\mathbf{p}_{\text{coarse}}$ and the essential-matrix and MAP solutions of both matching passes, collectively called the hypothesis set $\mathcal{P}$. To keep the pipeline simple, we select the guiding hypothesis as the fused estimate $\mathbf{p}_{\text{coarse}}$ and imagine a single view per pair.
% \noindent\textbf{位姿假设。} 粗估计阶段给出的位姿估计——融合估计 $\mathbf{p}_{\text{coarse}}$ 以及两遍匹配的本质矩阵解和 MAP 解——合称假设集 $\mathcal{P}$。为保持流程简单，选融合估计 $\mathbf{p}_{\text{coarse}}$ 作为引导假设，每个图像对只想象一个视图。

\noindent\textbf{Geometric Imagination.}
We observe that the scene seen from flight altitude is mostly planar and that the parallax of elevated structures such as rooftops is small relative to the image motion induced by the pose. Let $\tilde{\mathbf{x}}$ denote the homogeneous coordinates of a pixel in the source view, and $\mathbf{n}=(0,\sin\phi_S,\cos\phi_S)^\top$ denote the downward unit normal of the ground plane in the source camera frame. We can obtain its projected coordinates onto the target view $\tilde{\mathbf{x}}'$ by
\begin{equation}
\tilde{\mathbf{x}}'\sim\mathbf{H}\,\tilde{\mathbf{x}},\qquad
\mathbf{H}=\mathbf{K}\Bigl(\mathbf{R}+\frac{\mathbf{t}\,\mathbf{n}^{\top}}{h}\Bigr)\mathbf{K}^{-1},
\label{eq:homography}
\end{equation}
where $\mathbf{t}=-\mathbf{R}\mathbf{A}(\phi_S)^\top(\hat\ell_x,\hat\ell_y,\Delta\hat z)^\top$ is the metric translation given the guiding hypothesis and $h$ is the height of the source camera above the ground. The height of the camera can be estimated as the distance to the triangulated ground in the two-view reconstruction of the coarse stage, scaled to metres by the fused distance $\hat d$. When this reconstruction is not sufficiently reliable, as in pairs with few inliers, we replace $h$ by the median camera height $H_0$ of the training set. As indicated in Eq.~\eqref{eq:homography}, a key component of our refinement stage is a plane-induced renderer that imagines the source view $\widehat{S}$ by sampling pixels from the target view $G$ based on the homography $\mathbf{H}$, i.e., $\widehat{S}(\mathbf{x})=G(\mathbf{H}\mathbf{x})$. Note that the above formulation implicitly assumes 1) the ground is planar; 2) the cameras have zero roll and known intrinsics; 3) the guiding hypothesis is accurate enough for the two views to overlap after warping. If the first assumption is violated by elevated structures, the imagined view is locally displaced by a residual parallax, which the subsequent matching step tolerates.
% \noindent\textbf{几何想象。} 从飞行高度看，场景近似平面，屋顶等高出地面的结构所产生的视差相对于位姿引起的图像运动很小。设 $\tilde{\mathbf{x}}$ 为源视图像素的齐次坐标，$\mathbf{n}$ 为地平面在源相机坐标系中朝下的单位法向量，则它在目标视图上的投影由地面诱导的单应（式~\eqref{eq:homography}）给出，其中 $\mathbf{t}$ 是引导假设下的度量平移，$h$ 是源相机的离地高度。相机高度取粗估计阶段两视图重建中到三角化地面的距离，并用融合距离 $\hat d$ 换算为米；重建不够可靠时（如内点很少的图像对），用训练集的相机高度中位数 $H_0$ 代替。于是精修阶段的关键部件是一个平面诱导的渲染器：按 $\mathbf{H}$ 从目标视图 $G$ 采样像素，想象出源视图 $\widehat{S}$。该表述隐含三个假设：地面是平面；相机无横滚且内参已知；引导假设足够准确，使两视图在变换后有重叠。高出地面的结构违反第一个假设时，想象视图会在局部产生残余视差，由后续的匹配步骤容忍。

\noindent\textbf{Tile-Wise Re-Matching at Native Resolution.}
Because $\widehat{S}$ and $S$ are aligned up to a small residual, corresponding content falls into the same image window. We can leverage this property to match the overlap tile by tile at the native resolution, which is not possible before alignment. The frozen matcher then processes each of up to $3{\times}3$ co-located tiles with its full keypoint budget, yielding correspondences $(\mathbf{x}_i,\hat{\mathbf{x}}_i)$ between $S$ and $\widehat{S}$. Mapping every $\hat{\mathbf{x}}_i$ back to $G$ through $\mathbf{H}$ and gathering all tiles produces the imagined match set $\mathcal{M}_{\text{im}}=\{(\mathbf{x}_i,\mathbf{H}\hat{\mathbf{x}}_i)\}$. The relative pose can then be re-estimated from $\mathcal{M}_{\text{im}}$ with the same essential-matrix and MAP solvers. The residual displacement $\delta_i=\lVert\hat{\mathbf{x}}_i-\mathbf{x}_i\rVert_2$ reflects how far the imagined view is from the real observation at each match.
% \noindent\textbf{原始分辨率下的分块重匹配。} 由于 $\widehat{S}$ 与 $S$ 已对齐到只剩很小的残差，对应内容落在同一图像窗口内；利用这一点可以在原始分辨率上对重叠区逐块匹配（对齐之前做不到）。冻结的匹配器对最多 $3\times3$ 个位置对应的图块各用完整的关键点预算处理，得到 $S$ 与 $\widehat{S}$ 之间的对应点 $(\mathbf{x}_i,\hat{\mathbf{x}}_i)$。把每个 $\hat{\mathbf{x}}_i$ 经 $\mathbf{H}$ 映射回 $G$ 并汇总所有图块，得到想象匹配集 $\mathcal{M}_{\text{im}}$；再用同样的本质矩阵和 MAP 求解器由 $\mathcal{M}_{\text{im}}$ 重新估计相对位姿。残余位移 $\delta_i$ 反映了想象视图在每个匹配点处与真实观测相差多远。

\noindent\textbf{Hypothesis Verification.}
Assuming we have the hypothesis set $\mathcal{P}$ of the coarse stage, we can use $\mathcal{M}_{\text{im}}$ to verify each of its members. Each score is computed by testing $\mathcal{M}_{\text{im}}$ against the epipolar geometry of the corresponding hypothesis as
\begin{equation}
v_\tau(\mathbf{p}_k)=\frac{1}{|\mathcal{M}_{\text{im}}|}\sum_{i\in\mathcal{M}_{\text{im}}}\mathbf{1}\bigl[e_i(\mathbf{p}_k)<\tau\bigr],
\label{eq:verify}
\end{equation}
where $e_i(\mathbf{p}_k)$ is the Sampson error of match $i$ under hypothesis $\mathbf{p}_k$, and $\tau$ is a threshold controlling the strictness of the test. Instead of directly selecting the hypothesis with the best score, we pass the scores $v_\tau$ and the statistics of $\delta_i$ to the fusion as reliability features that complement the candidates themselves.
% \noindent\textbf{假设验证。} 有了粗估计阶段的假设集 $\mathcal{P}$，就可以用 $\mathcal{M}_{\text{im}}$ 验证其中每个假设：把 $\mathcal{M}_{\text{im}}$ 拿到该假设的对极几何下检验（式~\eqref{eq:verify}），$e_i(\mathbf{p}_k)$ 是匹配 $i$ 在假设 $\mathbf{p}_k$ 下的 Sampson 误差，$\tau$ 是控制检验严格程度的阈值。我们不直接选得分最高的假设，而是把得分 $v_\tau$ 和 $\delta_i$ 的统计量作为可靠性特征送入融合，与候选本身互补。

\noindent\textbf{Refined Fusion.}
We use the same reliability-weighted fusion for the coarse and refined stages, requiring a single additional set of regressors. It is the combination of the coarse estimate $\hat y$ and the two candidates re-estimated from $\mathcal{M}_{\text{im}}$:
\begin{equation}
y^\star=\operatorname{blend}\bigl(\hat y,\;c^{\text{im}}_{\mathrm{E}},\;c^{\text{im}}_{\text{MAP}}\bigr)
+r\bigl(\mathbf{f},\{v_\tau\},\{\delta_i\}\bigr),
\label{eq:refined}
\end{equation}
where $\operatorname{blend}(\cdot)$ denotes the reliability-weighted average of Eq.~\eqref{eq:blend}. For the distance, we use the same scale cues as in the coarse stage, recomputed from $\mathcal{M}_{\text{im}}$. In training, we leverage the held-out predictions of the coarse stage to match the test and training conditions and prevent the refinement from over-trusting its initialization.
% \noindent\textbf{精修融合。} 粗估计和精修两个阶段使用同一种按可靠性加权的融合，只需额外再训练一组回归器。精修结果是粗估计 $\hat y$ 与由 $\mathcal{M}_{\text{im}}$ 重新估计的两个候选的组合（式~\eqref{eq:refined}），其中 $\operatorname{blend}(\cdot)$ 即式~\eqref{eq:blend} 的可靠性加权平均。距离使用与粗估计阶段相同的尺度线索，但由 $\mathcal{M}_{\text{im}}$ 重新计算。训练时使用粗估计阶段的留出（held-out）预测，使训练与测试条件一致，避免精修过度信任其初始化。

\noindent\textbf{Canonical-View Heading Verification.}
The imagination in Eq.~\eqref{eq:homography} refines the relative pose but does not constrain the absolute heading $\psi_S$. We extend this paradigm with a second imagined view, rectifying each image to nadir, thus now a change of camera heading amounts to an in-plane rotation of the image. Each image is rectified via a rotation-only homography computed from its estimated tilt such that its tilt is zero -- its principal axis is then vertical. Unlike the plane-induced homography of Eq.~\eqref{eq:homography}, it is also exact for any scene depth, which removes the tilt as a nuisance factor and provides an exact rotation augmentation. The heading is discretized into $K$ rotations sampled at regular intervals and their mirrored copies. This verification benefits from a simple implementation by rotating the nadir view $K$ times and summing, after undoing each rotation, the von Mises natural parameters predicted by a second heading network. The estimate of the absolute heading obtained from the nadir views is then combined with the estimate from the coarse stage by summing their natural parameters, resulting in a more accurate estimate $\psi_S^\star$ of the source heading. The final prediction consists of the refined relative pose and the displacement rotated to the world frame with Eq.~\eqref{eq:world}: $\mathbf{p}^\star=(\Delta x^\star,\Delta y^\star,\Delta z^\star,\Delta\phi^\star,\Delta\theta^\star)$.
% \noindent\textbf{规范视图上的航向验证。} 式~\eqref{eq:homography} 的想象精修了相对位姿，但不约束绝对航向 $\psi_S$。我们把这一范式扩展到第二种想象视图：把每张图像校正为正下视，这样相机航向的变化就等价于图像的面内旋转。校正使用由估计俯仰角计算的纯旋转单应，使俯仰为零、主轴竖直；与式~\eqref{eq:homography} 的平面诱导单应不同，它对任意场景深度都是精确的，既去掉了俯仰这一干扰因素，又提供了精确的旋转增广。航向被离散为 $K$ 个等间隔旋转及其镜像；实现上把正下视图旋转 $K$ 次，撤销各自的旋转后，把第二个航向网络预测的 von Mises 自然参数相加。由正下视图得到的绝对航向估计再与粗估计阶段的估计通过自然参数相加合并，得到更准确的源相机航向 $\psi_S^\star$。最终预测由精修后的相对位姿和按式~\eqref{eq:world} 旋转到世界系的位移组成。

\noindent\textbf{Short-Range Mode for Closed-Loop Flight.}\label{sec:servo}
Estimating the translation near the goal is very challenging due to the lack of parallax, the degeneracy of the essential matrix, the bias of the fused distance at such baselines, and the scarcity of short baselines in the training pairs. As we show in the following, the imagined matches can effectively constrain the translation to overcome these challenges. We then obtain the constraint on the scaled translation $\mathbf{w}=\mathbf{t}/h$ for each imagined match $i$ as
\begin{equation}
\bar{\mathbf{x}}'_i\sim\mathbf{R}\,\bar{\mathbf{x}}_i+\mathbf{w}\,\bigl(\mathbf{n}^{\top}\bar{\mathbf{x}}_i\bigr),
\label{eq:servo}
\end{equation}
where $\bar{\mathbf{x}}=\mathbf{K}^{-1}\tilde{\mathbf{x}}$ denotes the normalized coordinates, with the rotation $\mathbf{R}$ fixed to its refined estimate. This is equivalent to solving a linear least-squares problem in $\mathbf{w}$, which remains well conditioned at zero baseline. Translation estimation is performed with the Direct Linear Transformation~\citep{hartley2003multiple} (DLT), which has a direct least-squares solution. The robust solution of Eq.~\eqref{eq:servo} is computed with iteratively reweighted least squares~\citep{holland1977robust}. This ensures robustness to mismatches and accounts for the residual parallax of elevated structures. To prevent this estimate from replacing the fusion where the latter is reliable, we apply it at a given step only if the refined distance is below 5\,m and at least 200 matches agree with the plane. The short-range mode does not require any additional model but only the camera height $h$ of Eq.~\eqref{eq:homography}, which converts $\mathbf{w}$ into a metric translation.
% \noindent\textbf{闭环飞行的近距离模式。} 在目标附近估计平移非常困难：视差不足、本质矩阵退化、融合距离在这种基线下有偏，而且训练图像对里短基线很少。下面说明，想象匹配可以有效地约束平移来克服这些困难。保持旋转 $\mathbf{R}$ 为精修后的估计，对每个想象匹配 $i$ 可以得到关于“按相机高度归一化的平移” $\mathbf{w}=\mathbf{t}/h$ 的约束（式~\eqref{eq:servo}），其中 $\bar{\mathbf{x}}=\mathbf{K}^{-1}\tilde{\mathbf{x}}$ 是归一化坐标。这等价于求解一个关于 $\mathbf{w}$ 的线性最小二乘问题，在零基线时仍然是良态的。平移用直接线性变换（DLT）求解，它有直接的最小二乘解；式~\eqref{eq:servo} 的鲁棒解用迭代重加权最小二乘计算，这样既对误匹配鲁棒，也能容忍高出地面的结构带来的残余视差。为了避免在融合仍然可靠的地方用它取代融合，只有当精修距离小于 5\,m 且至少有 200 个匹配与该平面一致时，才在这一步使用这个估计。近距离模式不需要任何额外的模型，只需要式~\eqref{eq:homography} 中的相机高度 $h$，用它把 $\mathbf{w}$ 换算成度量平移。
